% Canonical TeX edition. Initial text inherited from 3b067cc2eba4ac2f2f91248448727639d8ac7c5a
% Source: research/case_study/phase265_flag_relations.md
% Subsequent mathematical edits require a new audit.
% Options for packages loaded elsewhere
\PassOptionsToPackage{unicode}{hyperref}
\PassOptionsToPackage{hyphens}{url}
\documentclass[
  11pt,
  a4paper,
]{article}
\usepackage{xcolor}
\usepackage[margin=1in]{geometry}
\usepackage{amsmath,amssymb}
\setcounter{secnumdepth}{-\maxdimen} % remove section numbering
\usepackage{iftex}
\ifPDFTeX
  \usepackage[T1]{fontenc}
  \usepackage[utf8]{inputenc}
  \usepackage{textcomp} % provide euro and other symbols
\else % if luatex or xetex
  \usepackage{unicode-math} % this also loads fontspec
  \defaultfontfeatures{Scale=MatchLowercase}
  \defaultfontfeatures[\rmfamily]{Ligatures=TeX,Scale=1}
\fi
\usepackage{lmodern}
\ifPDFTeX\else
  % xetex/luatex font selection
  \setmainfont[]{STIXGeneral}
  \setmonofont[]{Menlo}
\fi
% Use upquote if available, for straight quotes in verbatim environments
\IfFileExists{upquote.sty}{\usepackage{upquote}}{}
\IfFileExists{microtype.sty}{% use microtype if available
  \usepackage[]{microtype}
  \UseMicrotypeSet[protrusion]{basicmath} % disable protrusion for tt fonts
}{}
\makeatletter
\@ifundefined{KOMAClassName}{% if non-KOMA class
  \IfFileExists{parskip.sty}{%
    \usepackage{parskip}
  }{% else
    \setlength{\parindent}{0pt}
    \setlength{\parskip}{6pt plus 2pt minus 1pt}}
}{% if KOMA class
  \KOMAoptions{parskip=half}}
\makeatother
\usepackage{longtable,booktabs,array}
\newcounter{none} % for unnumbered tables
\usepackage{calc} % for calculating minipage widths
% Correct order of tables after \paragraph or \subparagraph
\usepackage{etoolbox}
\makeatletter
\patchcmd\longtable{\par}{\if@noskipsec\mbox{}\fi\par}{}{}
\makeatother
% Allow footnotes in longtable head/foot
\IfFileExists{footnotehyper.sty}{\usepackage{footnotehyper}}{\usepackage{footnote}}
\makesavenoteenv{longtable}
\setlength{\emergencystretch}{3em} % prevent overfull lines
\providecommand{\tightlist}{%
  \setlength{\itemsep}{0pt}\setlength{\parskip}{0pt}}
\usepackage{newunicodechar}
\newfontfamily\numeralfallback{STIX Two Text}
\newunicodechar{∎}{\ensuremath{\blacksquare}}
\usepackage{fvextra}
\DefineVerbatimEnvironment{Highlighting}{Verbatim}{breaklines,breakanywhere,breakanywheresymbolpre={},breakanywheresymbolpost={},breaksymbolleft={},breaksymbolright={},commandchars=\\\{\}}
\newunicodechar{⁰}{{\numeralfallback\char"2070\relax}}
\newunicodechar{₀}{{\numeralfallback\char"2080\relax}}
\newunicodechar{₁}{{\numeralfallback\char"2081\relax}}
\newunicodechar{₂}{{\numeralfallback\char"2082\relax}}
\newunicodechar{₃}{{\numeralfallback\char"2083\relax}}
\newunicodechar{₄}{{\numeralfallback\char"2084\relax}}

\usepackage{seqsplit}
\renewcommand{\texttt}[1]{{\ttfamily\seqsplit{#1}}}
\setlength{\emergencystretch}{6em}

\newunicodechar{⁻}{{\numeralfallback\char"207B\relax}}
\newunicodechar{⁸}{{\numeralfallback\char"2078\relax}}
\newunicodechar{ⁱ}{{\numeralfallback\char"2071\relax}}
\newunicodechar{⁴}{{\numeralfallback\char"2074\relax}}
\newunicodechar{⁷}{{\numeralfallback\char"2077\relax}}
\newunicodechar{⁶}{{\numeralfallback\char"2076\relax}}
\newunicodechar{⁵}{{\numeralfallback\char"2075\relax}}
\newunicodechar{₈}{{\numeralfallback\char"2088\relax}}
\newunicodechar{ᵢ}{\textsubscript{i}}
\newunicodechar{ⱼ}{\textsubscript{j}}
\newunicodechar{ₜ}{\textsubscript{t}}
\newunicodechar{₅}{{\numeralfallback\char"2085\relax}}
\newunicodechar{₆}{{\numeralfallback\char"2086\relax}}
\newunicodechar{₇}{{\numeralfallback\char"2087\relax}}
\newunicodechar{ᵉ}{\textsuperscript{e}}
\newunicodechar{ᶠ}{\textsuperscript{f}}
\usepackage{bookmark}
\IfFileExists{xurl.sty}{\usepackage{xurl}}{} % add URL line breaks if available
\urlstyle{same}
% fallback for those not using the hyperref driver hyperxmp:
\makeatletter
\@ifundefined{xmpquote}{\newcommand{\xmpquote}[1]{#1}}{}
\makeatother
\hypersetup{
  pdftitle={Flag relations and derived incidence geometries},
  pdfauthor={Human Compression project},
  hidelinks,
  pdfcreator={LaTeX via pandoc}}

\title{Flag relations and derived incidence geometries}
\author{Human Compression project}
\date{11 October 2026}

\begin{document}
\maketitle
\section*{Plan of the common theory}
The two constructions studied here have different objects but share a practical route to their classification: describe chambers explicitly, identify every residue, and then check the remaining group action on a fiber. Part I retains the complete flag-relation theory for triangle geometries. Part II proves the corresponding residue and chamber statements for Steiner designs. Part III gives the complete unitary transport argument used by both applications.

The constructions and classification targets retain their original attribution. This edition assembles already accepted proofs, with their corrections; it asserts no new priority or additional source result. Each application supplies its own field, design and group hypotheses.
\part*{I. Flag relations and reconstruction}



\section{Scope and attribution}\label{scope-and-attribution}

This paper supplies a reusable interface for triangle constructions on
finite linear spaces. The construction is due to Delaby, Leemans and
Tranchida {[}DLT, Construction 1{]}. The common-neighbor reconstruction,
correlation extension and deleted-pencil description below reuse the
project's earlier \emph{Incidence compression}, §4. We give
complete proofs so that the present paper is an actual dependency of an
application, rather than a collection of references to an old combined
manuscript. The point-fiber synchronization in Theorem 2 also covers
two-point lines. The determinant criterion in Theorem 5 makes explicit
the semilinear alternative omitted from DLT's affine-plane
classification. Neither construction nor classical linear algebra is
claimed as a new subject. The accompanying application identifies
precisely which source statements these arguments recover or correct.

All sets in the geometric part are finite. A linear space S has a point
set P, a line set L, at least two points on each line, at least two
lines through each point, and a unique line through any two distinct
points. A flag is an incident pair (p,L). A uniform space has k points
on each line and r lines through each point; consequently v−1=r(k−1),
where v=\textbar P\textbar. We use ``automorphism'' for a
type-preserving incidence automorphism and ``correlation'' when
interchange of points and lines is allowed. A duality need not have
order two; a polarity does.

\section{1. A directed relation on
flags}\label{a-directed-relation-on-flags}

Write F for the flag set. Define the relation D by

\[
 (p,L)D(q,M)\quad\Longleftrightarrow\quad p\in M\text{ and }q\notin L.
\tag{1}
\]

Because q belongs to M, the conditions imply p≠q and M=pq≠L. If A and B
are the zero-one matrices for equality of point and equality of line on
F, then

\[
 D=(A-I)(B-I).
\tag{2}
\]

Indeed the only possible intermediate flag is (p,M). This observation
also proves that the integer product has entries zero or one.

The triangle incidence system Δ(S) has elements (f,i), where f∈F and
i∈Z/3Z. Distinct types i,i+1 are incident when fDg, respectively; the
reverse incidence is determined by symmetry. Its adjacency matrix, in
cyclic type order, is

\[
 \begin{pmatrix}0&D&D^{\mathsf T}\\D^{\mathsf T}&0&D\\
 D&D^{\mathsf T}&0\end{pmatrix}.
\tag{3}
\]

The cyclic shift τ(f,i)=(f,i+1) is therefore a triality. If p0,p1,p2 are
an ordered noncollinear triple, its chamber has flag of type i equal to
(pi,p(i−1)pi), with subscripts read cyclically. Conversely, (1) forces
these three joining lines in any chamber and excludes collinearity. A
panel between (p,L,0) and (q,pq,1) has exactly one completion for each
a∈L\{p\}: the third flag is (a,qa,2). Thus every flag extends to a
chamber, and every panel has k−1 completions in the uniform case. In
particular Δ(S) is firm precisely when k≥3 and is thin precisely when
k=3.

For a general rank-two input, the existence of a chamber of the literal
triangle system already forces a six-cycle in the input incidence graph:
its three point labels and three line labels are distinct, consecutive
incidences come from the chamber, and the required nonincidences
prohibit repeated labels. Consequently, if that nonempty triangle system
is a geometry, the input has gonality at most three. This recovers DLT
Proposition 3.3; it does not assert that every input of such gonality
gives a geometry.

The outgoing degree of the flag (p,L) is v−\textbar L\textbar, since
each q outside L gives the unique successor (q,pq). Hence flag
transitivity of Δ(S) forces every line of S to have the same size k: its
automorphism group is transitive on each vertex type. The identity
v−1=r\_p(k−1), obtained by partitioning P\{p\} among the lines through
p, then forces a constant r\_p=r. This proves uniformity without
assuming that the automorphisms of Δ(S) have already been reconstructed.

\section{2. Reconstruction from common
neighbors}\label{reconstruction-from-common-neighbors}

\textbf{Lemma 1 (point and line signatures).} For distinct flags f,g in
a uniform linear space, the following alternatives hold:

{\def\LTcaptype{none} % do not increment counter
\begin{longtable}[]{@{}lll@{}}
\toprule\noalign{}
Relation between f and g & Common successors & Common predecessors \\
\midrule\noalign{}
\endhead
\bottomrule\noalign{}
\endlastfoot
same point & (r−2)(k−1) & not needed \\
different points & 0 or k−2 & not needed \\
same line & not needed & (k−2)(r−1) \\
different lines & not needed & 0 or r−2 \\
\end{longtable}
}

For two flags (p,L),(p,M) at the same point, a common successor uses a
line through p different from L and M, and then one of its k−1 other
points. For flags at different points p,q, its line must be pq. If
neither forbidden line is pq, its point may be any of the k−2 points on
pq other than p and q; otherwise there are no choices. Reversing the
reasoning proves the last two rows. Distinct lines either meet in one
point or do not meet, so the reverse count has no further possibilities.
This proves the lemma.

\textbf{Theorem 2 (type-preserving reconstruction).} If r≥3 and k≥2,
restriction to the underlying flags identifies Aut(Δ(S)) with Aut(S).

The value (r−2)(k−1) is strictly larger than k−2. Thus the maximum
off-diagonal common-successor relation on each type identifies precisely
which flags have the same point. Every type-preserving automorphism φ of
Δ therefore induces a permutation fi of P on type i. Write

\[
 \phi(p,L,i)=(f_i(p),M_i(p,L),i).
\]

The point labels occurring among the successors of (p,L,i) are exactly
P\textbackslash L. Applying φ gives

\[
 P\setminus M_i(p,L)=f_{i+1}(P\setminus L),
 \qquad M_i(p,L)=f_{i+1}(L).
\tag{4}
\]

Since the image is a flag, fi(p) belongs to every f(i+1)(L) with L
through p. At least two distinct lines pass through p and their
intersection is \{p\}. Intersecting these images shows fi(p)=f(i+1)(p).
All three point permutations are therefore a single permutation f, and
(4) shows that f maps every line to a line and that φ is its diagonal
action on flags. The inverse satisfies the same condition, so f is an
automorphism of S. Conversely every automorphism of S preserves (1),
proving the claimed equality. The same argument, with the target point
set written P′, applies to any isomorphism Δ(S)→Δ(S′) between uniform
spaces with the same parameters k,r: equation(4) becomes
P′\textbackslash M\_i=f\_(i+1)(P\textbackslash L), and the same
intersections synchronize the three point bijections. Thus every such
triangle isomorphism is induced by one unique isomorphism S→S′. Notice
that the proof never needs k−2\textgreater0; two-point lines are
included.

If k,r≥3, both exceptional common-neighbor values in Lemma 1 are strict
maxima. A correlation reversing the cyclic type order exchanges the
successor and predecessor signatures and therefore exchanges the two
flag partitions. On each type it therefore induces a bijection from
points to lines and from lines to points. A point fiber and a line fiber
meet in a flag exactly when their labels are incident, so these two
bijections form a duality δi of S. In particular the action on type 0 is
induced by one duality δ0. The globally defined lift of δ0 reverses the
cyclic type order, since it transposes D. Composing the given
correlation with its inverse lift gives a type-preserving automorphism
which is the identity on type 0. Theorem 2 forces that automorphism to
be the identity on every type. Thus all three actions were induced by
δ0. No extra pointwise kernel has been suppressed in this
reconstruction.

\textbf{Theorem 3 (the complete correlation extension).} For k,r≥3 there
is a natural isomorphism

\[
 \operatorname{Cor}(\Delta(S))\simeq C_3\rtimes_\epsilon
 \operatorname{Cor}(S),
\tag{5}
\]

where ε records whether a correlation preserves or interchanges the two
types of S, and ε=−1 acts on C3 by inversion. If S has no duality, (5)
is C3×Aut(S). If S has a duality, the different exact sequence with
kernel Aut(S) and quotient S3 splits if and only if S has a polarity.

For completeness, a correlation δ of S acts on a flag by δ(p,L)=(δp,δL)
when it preserves types and by (δL,δp) when it interchanges them. Define

\[
 \widetilde\delta(f,i)=(\delta f,\epsilon(\delta)i).
\]

This is a homomorphism into Cor(Δ), fixes type 0, and conjugates τ to τ
raised to ε(δ). The reconstruction just proved shows that τ and all
these lifts generate the full correlation group. The two displayed
factors have trivial intersection, proving (5). A polarity supplies an
involutory lift of a type transposition and, together with τ, a
complement S3. Conversely any complement has such an involutory lift,
whose induced duality of S has order two. This proves the splitting
criterion. A polarity is not required for (5) itself; that qualification
belongs only to the latter exact sequence.

The chamber correspondence and Theorem 2 give the useful equivalence

\[
 \Delta(S)\text{ is flag-transitive}\quad\Longleftrightarrow\quad
 \operatorname{Aut}(S)\text{ is transitive on ordered noncollinear triples}.
\tag{6}
\]

The same assertion holds for any specified subgroup of Aut(S), using its
induced diagonal action.

\section{3. Connectivity and the deleted-pencil
residue}\label{connectivity-and-the-deleted-pencil-residue}

For k,r≥3, Δ(S) is connected. In a fixed type, flags with the same point
have a common successor and flags with the same line have a common
predecessor, by Lemma 1. These moves connect all flags because the
incidence graph of a linear space is connected: two points have their
joining line. Thus each type lies in one component, and an edge between
successive types joins the three components.

\textbf{Theorem 4 (residue interface).} The residue at (p,L,0) is the
bipartite incidence structure with parts

\[
 \{Q:Q\notin L\},\qquad
 \{(a,M):a\in L\setminus\{p\},\ M\ne L,\ a\in M\},
\tag{7}
\]

and incidence Q∈M. For the first part, the corresponding original vertex
is (Q,pQ,1); the second part already gives the original vertex of type
2. The residue in the Van Nypelseer construction used in L092 has
precisely the same description with a allowed to range over all of L.
Formula (7) is obtained by deleting one pencil, not by assuming that
this deletion preserves connectivity.

To verify (7), incidence with (p,L,0) forces a type-1 vertex to be
(Q,pQ,1) for Q outside L, and a type-2 vertex to be (a,M,2) with
a∈L\{p\} and M≠L. Their mutual incidence requires Q∈M and a∉pQ. The
latter condition is automatic, since a≠p and pQ meets L only in p.~This
proves both directions and the claimed literal identification.

In projective or affine dimension at least three, each edge in (7) stays
inside the plane generated by L and Q. Different such planes have
disjoint outside-point sets and disjoint allowed-line sets. Thus the
residue is a disjoint union of its plane residues. Whenever those plane
residues are connected, this is its exact component decomposition.

\section{4. Determinants classify affine-plane frame
orbits}\label{determinants-classify-affine-plane-frame-orbits}

We write vectors as rows. Let K be a finite field and H a subgroup
between SL(d,K) and ΓL(d,K), with d≥2. A semilinear element (A,σ) acts
as v↦v\^{}σ A. For the affine group V⋊H, translating the first point of
an ordered noncollinear triple to zero identifies its orbit with the
H-orbit of an ordered independent pair of vectors.

\textbf{Theorem 5 (frame-orbit criterion).} If d≥3, H is transitive on
ordered independent pairs. If d=2, its orbits on ordered bases are in
bijection with the orbits on K* of its induced determinant-semilinear
transformations

\[
 t\longmapsto \det(A)\,\sigma(t).
\tag{8}
\]

In dimension two, the affine group is transitive on ordered noncollinear
triples exactly when the action (8) is transitive. For a linear
H≤GL(2,K) this is equivalent to H=GL(2,K). For a semilinear H it need
not imply containment of GL(2,K).

For d≥3, extend either ordered independent pair to a basis. A linear map
sending one complete basis to the other may have determinant c≠1.
Multiply one of the target basis vectors outside the given pair by c⁻¹.
The resulting map has determinant one and still sends the required pair
correctly.

For d=2, arrange an ordered basis as the two rows of an invertible
matrix B. If B and C have equal determinant, there is exactly one matrix
M=B⁻¹C with BM=C, and det M=1. Thus SL(2,K) acts regularly on every
determinant fiber. A semilinear transformation sends B to B\^{}σ A and
its determinant to (8). These facts prove both implications of the orbit
correspondence: an H-orbit has one orbit of determinant values;
conversely first move a determinant to the desired value and then use
the contained SL group within that fiber. The transformations (8) form a
group because they are induced by an actual action on matrices. For a
linear H, the determinant image is a subgroup of K* acting by
multiplication. Transitivity forces the full image K*; together with its
contained kernel SL(2,K), this gives H=GL(2,K).

The statements about equal-determinant fibers, semilinear determinant
transport, and orbit saturation are formalized in
\texttt{formal/DeterminantFibers.lean}. The geometric
classification, finite-field subgroup identifications and connectivity
proofs are written mathematics, not consequences already formalized in
that file.

\section{5. A classical projective-line normalizer
interface}\label{a-classical-projective-line-normalizer-interface}

We record an elementary proof of the normalizer used by the circle
branch of the application, so that it is not hidden in a
higher-dimensional collineation theorem. This is a classical result, not
a novelty claim.

\textbf{Theorem 6.} For the natural action on P¹(K), where K is any
finite field,

\[
 N_{\operatorname{Sym}(\mathbf P^1(K))}
   (\operatorname{PSL}(2,K))=\operatorname{P\Gamma L}(2,K).
\tag{9}
\]

Let f normalize PSL. Since PGL is sharply three-transitive and itself
normalizes PSL, composing f with a projectivity lets us suppose that f
fixes ∞,0,1. The stabilizer of ∞ in PSL consists of x↦ax+b with a a
nonzero square (all nonzero a in even characteristic). Its translation
subgroup U=\{x↦x+b\} is its unique Sylow p-subgroup, since U is normal
and the scalar factor has order prime to p.~Thus f normalizes U. Because
f fixes0, conjugation sends the translation by b to the translation by
f(b), giving

\[
 f(x+b)=f(x)+f(b).
\tag{10}
\]

The stabilizer of ∞ and0 in PSL consists of the square multipliers.
Since f also fixes1, conjugating the multiplier by a square s gives the
multiplier by f(s), and hence f(sx)=f(s)f(x) for every nonzero square s;
the assertion also holds for s=0. In even characteristic every field
element is a square, so f is multiplicative. In odd characteristic every
t∈K is a sum of two squares: the set of squares including0 and its
translate t minus that set both have (q+1)/2 elements, so they
intersect. Write t=s+u with s,u squares. Additivity gives

\[
 f(tx)=f(sx)+f(ux)=(f(s)+f(u))f(x)=f(t)f(x).
\]

Thus f is a field automorphism in all cases. Undoing the initial
projectivity proves containment in PΓL. Conversely projectivities and
field automorphisms normalize PSL by their matrix action, proving
equality. The proof includes q=2,3: the same translation and square
arguments apply, without assuming simplicity of PSL. ∎

\section{References and dependency
boundary}\label{references-and-dependency-boundary}

{[}DLT{]} R. Delaby, D. Leemans and P. Tranchida, \emph{Geometries with
trialities arising from linear spaces}, Algebraic Combinatorics 9
(2026), 379--402, DOI10.5802/alco.482. Construction (1) and the
geometric targets are theirs.

{[}IC{]} \emph{Incidence compression}, repository manuscript, §4. Its
flag-matrix, correlation and deleted-pencil arguments are reused above.
The application must charge this complete rewritten dependency; citing
it does not make its proof cost disappear.

The group classification of locally two-transitive linear spaces remains
an explicitly retained application input. The complete Hermitian norm-shell
proof is included in Part III below. A
source correction, an independently verified new question answer, and a
recovered principal contract have different meanings and receive
separate accounting.


\part*{II. Derived designs and chamber fibers}
\section{The design construction}

Let $\mathcal S=S(t,k,v)$ be a finite Steiner system on a set $\Omega$:
every $t$-subset lies in exactly one block, and every block has $k$
points. We assume $2\leq t\leq k\leq v$ and write
\[
 r=\frac{v-t+1}{k-t+1}.
\]
This is the number of blocks through a fixed $(t-1)$-subset. Indeed,
their complements to that subset partition the remaining points into
sets of size $k-t+1$. Distinct blocks through that subset intersect
precisely in it. These elementary observations will account for the
exceptional incidences in the construction.

We use Construction 2 of Leemans [VN], with its attribution to Van
Nypelseer and Meixner. The incidence structure $\Gamma(\mathcal S)$ has
ordinary objects of types $1,\ldots,t-2$, namely the subsets of those
sizes. Its three remaining types are marked objects
\[
 m=(T,E),\qquad |T|=t-1,\quad T\cap E=\varnothing,
 \quad T\cup E\text{ a block},
\]
blocks $B_0$, and copied points $z\in\Omega$. Ordinary objects of
different types are incident by inclusion. Their incidences with the
last three types, and the incidences between those types, are
\[
\begin{array}{lll}
 X\sim(T,E)\iff X\subseteq T,&
 X\sim B_0\iff X\subseteq B_0,&
 X\sim z\iff z\notin X,\\
 (T,E)\sim B_0\iff B_0\cap(T\cup E)=T,&
 (T,E)\sim z\iff z\in E,&
 B_0\sim z\iff z\notin B_0.
\end{array}                                                    \tag{S1}
\]
Objects of the same type are incident only when equal. At $t=2$ the
ordinary part is empty. A marked object is then an original incident
point--line pair, so (S1) is exactly the rank-three Construction 3.

\paragraph{Theorem S.1 (chambers and firmness).}
For a nontrivial Steiner system with $t\geq2$, the construction is a
firm geometry if and only if
\[
 k\geq t+1,\qquad r\geq3.                                    \tag{S2}
\]
Under these conditions its panels missing an ordinary type have two
completions. The panels missing the marked, block and copied-point
types have respectively $k-t+2$, $r-1$ and $k-t+1$ completions.

\emph{Proof.}
The ordinary objects in a flag form a chain. Suppose first that a
marked object $(T,E)$ is present. Complete that chain inside $T$, choose
an alternative block through $T$ if none is present, and choose a copied
point of $E$ if necessary. Distinct blocks through $T$ have disjoint
tails, so these choices respect any already specified exceptional
objects and give a chamber.

If a block $B_0$ is present but no marked object is specified, extend
the largest ordinary set to a $(t-1)$-subset $T$ of $B_0$. Use the
prescribed copied point, or choose $z\notin B_0$ if none is prescribed.
The latter is possible since $r\geq3$ implies $v>k$. The unique block
$B'$ through $T\cup\{z\}$ differs from $B_0$ and meets it in $T$.
The marked object $(T,B'\setminus T)$ completes the required data.
If neither a block nor a marked object is present, use or choose a
copied point $z$ outside the largest ordinary set, extend that set to
$T\subseteq\Omega\setminus\{z\}$, and take the block on
$T\cup\{z\}$ and any alternative block through $T$. Complete the
chain inside $T$ in each case. This also treats the empty ordinary
chain when $t=2$, and proves that every flag extends to a chamber.

For a missing ordinary type the bounding subsets differ by two
elements, giving exactly two intermediate subsets. For a missing
marked type let $R$ be the largest ordinary set, of size $t-2$; put
$R=\varnothing$ if $t=2$. Choose $T=R\cup\{u\}$ with
$u\in B_0\setminus R$. The unique block through $T\cup\{z\}$
then determines the marked object, giving $k-t+2$ choices. The other
two panel counts are the number of alternative blocks through $T$
and the number of points of $E$. They are at least two under (S2).

Conversely, in a firm geometry a marked object belongs to a chamber.
Removing the copied point or the alternative block gives panels with
$k-t+1$ and $r-1$ completions. Both must have size at least two.
If $r=1$, no chamber contains that marked object in the first place.
This proves necessity as well. \hfill$\square$

In the remainder of this part we assume (S2).

\paragraph{Lemma S.2 (the induced automorphism group).}
Under (S2), the natural action identifies
$\operatorname{Aut}(\Gamma(\mathcal S))$ with
$\operatorname{Aut}(\mathcal S)$, where automorphisms preserve every
type.

\emph{Proof.}
An automorphism permutes the copied points. Their nonincidences with
blocks recover exactly each block as a subset of $\Omega$, so the
induced permutation is a design automorphism. Nonincidences with an
ordinary object recover its point subset too. For a marked object,
its copied-point neighbors recover $E$. Its block neighbors are the
$r-1\geq2$ alternative blocks through $T$; the intersection of any
two of them is $T$. Thus $T$ is recovered as well. All objects are
therefore determined by this one point permutation. Conversely a
design automorphism preserves all six rules (S1), proving the claim.
\hfill$\square$

\section{Chambers as one stabilizer fiber}

For a group acting equivariantly on a surjection $Y\to X$ between
nonempty sets, transitivity
on $Y$ is equivalent to transitivity on $X$ and transitivity of one
point stabilizer on the corresponding fiber. To prove the converse,
first move the image of a given element to the target image, and then
move within that fiber. For the forward implication, an element carrying
two points in the same fiber necessarily fixes their common image.
We include this standard action argument to make its precise use clear.

\paragraph{Theorem S.3 (the chamber fiber).}
For every $G\leq\operatorname{Aut}(\mathcal S)$ there is a
$G$-equivariant bijection
\[
\operatorname{Ch}(\Gamma(\mathcal S))\longleftrightarrow
\left\{((x_1,\ldots,x_t),B_0):
\begin{array}{l}
x_1,\ldots,x_t\text{ are distinct},\\
\{x_1,\ldots,x_{t-1}\}\subset B_0,\quad x_t\notin B_0
\end{array}\right\}.                                        \tag{S3}
\]
The fiber over each ordered $t$-tuple has size $r-1$. Consequently
$G$ is chamber-transitive if and only if it is $t$-transitive on
$\Omega$ and the stabilizer of one ordered $t$-tuple is transitive
on those $r-1$ blocks. In particular the chamber count is
\[
 (v)_t(r-1),\qquad (v)_t=v(v-1)\cdots(v-t+1).                 \tag{S4}
\]

\emph{Proof.}
The full ordinary chain and its containing $(t-1)$-set $T$ order the
points of $T$: successive differences are singletons. Read these as
$x_1,\ldots,x_{t-1}$ and take $x_t$ to be the copied point. The
alternative block is $B_0$. Conversely the data on the right of (S3)
determine the unique block $B'$ on
$T\cup\{x_t\}$. It differs from $B_0$, and the Steiner property
gives $B'\cap B_0=T$. The marked object is
$(T,B'\setminus T)$ and the ordinary chain consists of the initial
subsets of the tuple. The two constructions are inverse and commute
with the group action. Of the $r$ blocks through $T$, exactly one
contains $x_t$. The preceding stabilizer argument and elementary
counting now give all assertions. \hfill$\square$

This also proves the exact formulation of Property $(*)$ in
[VN, Theorem 9]. The source's Construction 1 has ordinary subsets of
sizes $1,\ldots,t-1$ and blocks, incident by inclusion. A chamber of
our geometry is equivalently a chamber of Construction 1, an alternative
block $B'$ through its final $(t-1)$-set $T$, and a point
$z\in B'\setminus T$. Forgetting $z$ and then $B'$ gives two
surjective equivariant maps. Applying the stabilizer argument at each
map says precisely that the action on Construction 1 is transitive,
its chamber stabilizer is transitive on the alternative blocks, and
the stabilizer of an alternative block is transitive on its remaining
points. This is Property $(*)$. The alternative encoding (S3) proves
that these conditions imply $t$-transitivity. Transitivity on the first
choice alone would not suffice. The numbers of objects are
\[
 n_i=\binom vi\quad(1\leq i\leq t-2),\qquad
 n_{\rm block}=\frac{\binom vt}{\binom kt},\qquad
 n_{\rm marked}=n_{\rm block}\binom{k}{t-1},\qquad
 n_{\rm copy}=v.                                             \tag{S5}
\]
The marked count chooses its block and then $T$. These formulas and
the panel counts determine numerical labels, but do not identify a
stabilizer group up to isomorphism.

\section{Derivation and all residues}

For $X\subseteq\Omega$ of size $j<t$, write $\mathcal S_X$ for
the derived system: its points are $\Omega\setminus X$, and its
blocks are $B\setminus X$ for blocks containing $X$. Its parameters
are $S(t-j,k-j,v-j)$. In particular $k-t+1$ and $r$ are unchanged.
Let $\mathcal B(X)$ be the geometry of the nonempty proper subsets
of a finite set $X$, with inclusion incidence and one type per size.
A join makes every object of one factor incident with every object
of another. Empty rank-zero factors are allowed.

\paragraph{Lemma S.4 (Boolean residues and the derived residue).}
The geometry $\mathcal B(X)$ and its truncation to all subset sizes
$1,\ldots,h<|X|$ are residually connected. Under (S2), the residue
at an ordinary object $X$ of size $j$ is naturally
\[
 \Gamma(\mathcal S)_X\simeq
 \mathcal B(X)*\Gamma(\mathcal S_X).                          \tag{S6}
\]

\emph{Proof.}
A flag in a subset geometry is a chain. Its residue is the join of
subset geometries on successive set differences, with a final
truncated factor above its largest member. In a factor of rank at
least two, singleton objects are connected through two-element
objects, and every other object has a singleton neighbor. If two
positive-rank factors remain, their complete cross-incidence connects
the whole residue even when each factor has rank one. Rank-zero and
rank-one residues impose no connectivity condition. This proves the
first assertion and also proves that joins of residually connected
geometries are residually connected.

Objects below $X$ give $\mathcal B(X)$. Above $X$, send an ordinary
$Y$ to $Y\setminus X$, a marked $(T,E)$ to $(T\setminus X,E)$,
a block $B_0$ to $B_0\setminus X$, and an incident copied point
$z\notin X$ to itself. These are bijections onto the objects of
$\Gamma(\mathcal S_X)$. Subtracting $X$ preserves inclusions and
all the exceptional rules (S1). Every object below $X$ is incident
with every object of this second factor, including its copied points.
Thus (S6) is an incidence isomorphism, not just an equality of counts.
\hfill$\square$

\paragraph{Theorem S.5 (reduction to rank three).}
Under (S2), $\Gamma(\mathcal S)$ is residually connected if and
only if $\Gamma(\mathcal S_R)$ is residually connected for every
$(t-2)$-subset $R$. At $t=2$ this means the identity assertion with
$R=\varnothing$. No transitivity assumption is required.

\emph{Proof.}
For necessity take a complete ordinary chain with largest member
$R$, using every size from $1$ to $t-2$. Its residue is exactly
$\Gamma(\mathcal S_R)$: the complete chain removes the Boolean
factors in (S6). Every flag residue of a residually connected
geometry is residually connected. Connectedness of the join at $R$
alone would not prove this implication, so the complete chain matters.

For sufficiency induct on $t$. The starting case $t=2$ is tautological.
For an ordinary $X$ of size $j$, the final two-design derivatives of
$\mathcal S_X$ are precisely the $\mathcal S_R$ with
$R\supseteq X$, $|R|=t-2$. Induction and Lemma S.4 imply that the
residue at $X$ is residually connected. This deals with every flag
containing an ordinary object, including flags with additional
exceptional objects.

It remains to inspect the subsets of the three exceptional types.
At a marked object $(T,E)$ the residue is the join of
$\mathcal B(T)$, the $r-1$ alternative blocks through $T$, and
the discrete set $E$ of copied points. All cross-incidences are
present. Fixing an additional block or copied point removes its
discrete factor, so Lemma S.4 treats all four cases containing the
marked object whenever the remaining rank is at least two.

At an incident block--copied-point pair $(B_0,z)$, the marked objects
are in bijection with the $(t-1)$-subsets $T$ of $B_0$: use the unique
block through $T\cup\{z\}$. The remaining incidence is precisely
the subset geometry on $B_0$ with sizes $1,\ldots,t-1$, which is
connected through points and two-subsets for $t\geq3$.

At a block $B_0$ alone, its ordinary singleton points and the copied
points outside it form a nonempty complete bipartite graph. Every
other ordinary or marked object has a singleton neighbor in $B_0$.
At a copied point $z$ alone, the ordinary and marked objects give
all subsets of $\Omega\setminus\{z\}$ of sizes $1,\ldots,t-1$;
the largest subsets again determine their unique blocks through $z$.
This connected inclusion geometry has a point neighbor on every
remaining block. Finally the full graph contains the graph between
ordinary singletons and copied points with adjacency $x\ne z$.
It is $K_{v,v}$ minus its matching, connected since $v\geq3$;
every other object has a singleton neighbor. These are the remaining
four cases. Together they exhaust the exceptional-only flags and
complete the induction. \hfill$\square$

\paragraph{Theorem S.6 (rank-two intersections).}
Under (S2), every rank-two residue of $\Gamma(\mathcal S)$ is a
partial linear space or a generalized digon. Thus the source's
property $(\mathrm{IP})_2$ holds throughout this family.

\emph{Proof.}
At $t=2$, the residue at an original flag $(x,M)$ is complete
bipartite between $M\setminus\{x\}$ and the other lines through
$x$. At an original point $p$, a remaining flag is determined by
its point $x$, since its line is $px$; incidence with the lines
avoiding $p$ is $x\in L$. This is a subincidence structure of the
original linear space. At an original line $L$, a remaining flag
$(x,M)$ is determined by $M$, since $x=M\cap L$; its incidence
with points outside $L$ is again a subincidence structure of the
original linear space. The latter two residues are partial linear
spaces.

For higher $t$, a defining flag containing an ordinary object is
handled by (S6). Residues of a join are joins of factor residues.
If two positive-rank factors remain they both have rank one and
give a generalized digon. If only the derived factor remains use
induction. In a Boolean factor, a rank-two residue is either a
join of two rank-one intervals or the incidence of singletons and
two-subsets of a three-element set. These give the same two possibilities.

A defining flag containing only exceptional objects has size $t-1$,
so here $t\leq4$. For $t=4$ all three exceptional objects are
fixed and the remaining objects are the nonempty proper subsets
of their three-set $T$, giving a triangle. For $t=3$, a marked
object and a block leave the two points of $T$ and the copied points
of $E$ with complete incidence. A marked object and a copied point
similarly leave the two points of $T$ and the alternative blocks.
A block $B_0$ and copied point $z$ leave the points and two-subsets
of $B_0$; each two-subset determines its marked object through $z$.
This is a partial linear space. These cases exhaust all possibilities.
\hfill$\square$

The Boolean argument uses the complete-chain interval decomposition.
It does not assert that arbitrary two-type truncations of a Boolean
lattice are partial linear spaces. The source constructions, firmness
criterion and transitivity target are retained from [VN]; the residue
and fiber proofs above reuse the accepted phase182/186 proof and its
independent phase187/199 reviews. The all-rank induction and its
applications are written mathematics, not a claimed full Lean theorem.

\paragraph{Reference.}
[VN] D. Leemans, \emph{A generalization of a construction due to Van
Nypelseer}, Beitr\"age zur Algebra und Geometrie 46 (2005), 561--574,
Constructions 2--3 and Theorems 1 and 9--11. The two retained
classification inputs used in the companion application are stated
there with their exact hypotheses; they do not follow from S.5.

\part*{III. Hermitian transport and residues}
\section{Transport on a nonzero Hermitian shell}

The difficult residues in both applications can be described on the
same set of vectors. We give the complete algebraic argument here,
including the maps that turn its group generators into actual incidence
paths. These are elementary unitary calculations, not a claim of a new
classification theorem. The proof combines the accepted unital argument
with its explicit transport supplement.

Let $K$ be a commutative field and $\sigma$ an involutive field
automorphism, written $a\mapsto\bar a$. On column vectors in $K^2$
put
\[
 h((x,y),(u,v))=x\bar v+y\bar u,
 \qquad N(x,y)=x\bar y+\bar x y,\qquad
 J=\begin{pmatrix}0&1\\1&0\end{pmatrix}.
\]
Write $\operatorname{SU}(h)$ for the determinant-one matrices
preserving this form. For $M=\left(\begin{smallmatrix}a&b\\c&d\end{smallmatrix}\right)$
with determinant one, the equation $M^{\mathsf T}J\sigma(M)=J$
is equivalent to
\[
 \bar a=a,\quad \bar d=d,\quad \bar b=-b,\quad \bar c=-c,
 \qquad ad-bc=1.                                             \tag{H1}
\]
Indeed solving for $\sigma(M)$ gives
$\sigma(M)=J(M^{\mathsf T})^{-1}J
=\left(\begin{smallmatrix}a&-b\\-c&d\end{smallmatrix}\right)$;
the calculation is reversible. Direct expansion also verifies
$N(Mv)=N(v)$ under (H1).

\paragraph{Theorem H.1 (unique unitary transport).}
If $N(v)=N(w)=k\ne0$, there is exactly one element of
$\operatorname{SU}(h)$ carrying $v$ to $w$. For $v=(x,y)$ and
$w=(u,r)$ it is
\[
 T_{w,v}=\frac1k
 \begin{pmatrix}
 u\bar y+\bar u y & u\bar x-\bar u x\\
 r\bar y-\bar r y & r\bar x+\bar r x
 \end{pmatrix}.                                              \tag{H2}
\]

\emph{Proof.}
The columns of
\[
 B_v=\begin{pmatrix}x&-\bar x\\y&\bar y\end{pmatrix}
\]
are orthogonal and have respective norms $k,-k$. Its determinant
is $k$, so it is invertible. Thus $B_wB_v^{-1}$ has determinant
one, preserves the Hermitian form, and carries the first column of
$B_v$ to that of $B_w$. Matrix multiplication gives (H2). The
argument includes characteristic two: the two columns remain
independent because their determinant is the nonzero element $k$.

For uniqueness, (H1) implies that any such matrix taking $(x,y)$
to $(u,r)$ must take $(-\bar x,\bar y)$ to
$(-\bar u,\bar r)$. Its values on the basis formed by the columns
of $B_v$ are determined, so it is $B_wB_v^{-1}$.
\hfill$\square$

\section{Two isotropic directions suffice}

For $\bar t=-t$, the matrices
\[
 U_t=\begin{pmatrix}1&t\\0&1\end{pmatrix},\qquad
 V_t=\begin{pmatrix}1&0\\t&1\end{pmatrix}                     \tag{H3}
\]
satisfy (H1). Their displacements lie on the two isotropic coordinate
axes. If $N(v)\ne0$ and $t\ne0$, those displacements are nonzero:
neither coordinate of $v$ can be zero.

Now let $K=\mathbb F_{q^2}$ with $\sigma(a)=a^q$. The trace-zero
space has dimension one over $\mathbb F_q$. Choose a nonzero
$\delta$ with $\bar\delta=-\delta$, including in characteristic
two. Then $\delta^2\in\mathbb F_q^*$, and the entries of (H1)
satisfy $a,d\in\mathbb F_q$, $b,c\in\delta\mathbb F_q$.
Conjugation by $\operatorname{diag}(\delta,1)$ identifies these
matrices with $\operatorname{SL}(2,q)$: explicitly, conjugating
by its inverse on the left changes the off-diagonal entries to
$b/\delta$ and $c\delta$, which range over $\mathbb F_q$.
The two families (H3) therefore become all upper and lower
elementary matrices over $\mathbb F_q$.

For clarity, these elementary matrices generate $\operatorname{SL}(2,F)$
over every field $F$. If
$A=\left(\begin{smallmatrix}a&b\\c&d\end{smallmatrix}\right)$
has $a\ne0$, left multiplication by $V_{-c/a}$ and right
multiplication by $U_{-b/a}$ give $\operatorname{diag}(a,a^{-1})$.
If $a=0$, then $c\ne0$ and $U_1A$ has nonzero upper-left entry.
Finally, for $t\ne0$,
\[
 W(t)=U_tV_{-t^{-1}}U_t=
 \begin{pmatrix}0&t\\-t^{-1}&0\end{pmatrix},\qquad
 W(t)W(-1)=\operatorname{diag}(t,t^{-1}).                     \tag{H4}
\]
This proves generation, including over $\mathbb F_2$.

\paragraph{Theorem H.2 (connectivity using specified directions).}
Let $W$ be a nondegenerate Hermitian plane over $\mathbb F_{q^2}$,
let $k\in\mathbb F_q^*$, and let $\mathcal A$ be a set of its
isotropic one-spaces containing at least two distinct members.
The graph on $X_k=\{v\in W:h(v,v)=k\}$ which joins distinct
$v,w$ whenever $K(v-w)\in\mathcal A$ is connected. The group
$\operatorname{SU}(W)$ acts regularly on $X_k$, and
$|X_k|=q(q^2-1)$.

\emph{Proof.}
Choose isotropic representatives $e,f$ for two specified directions.
Their pairing is nonzero, since the orthogonal complement of an
isotropic line in a nondegenerate plane is that same line. Rescale
them to give Gram matrix $J$. Theorem H.1 proves regularity of
the special unitary action on each nonempty nonzero shell. Such
shells are nonempty: with $y=1$ the equation $N(x,1)=k$ is the
surjective trace equation. The two root families in this chosen
basis generate the whole group by (H4), and each nonidentity
root element makes a nonzero move in one of the two allowed
directions. An element carrying one shell vector to another can
therefore be written as a product of moves along graph edges.
This proves connectivity. Finally $|\operatorname{SL}(2,q)|
=q(q^2-1)$, by choosing a nonzero first column and then one of
the $q$ second columns with determinant one. Regularity gives
the claimed shell size. \hfill$\square$

\section{The projective incidence map}

Let $V$ be a nondegenerate Hermitian three-space over
$K=\mathbb F_{q^2}$. Its classical unital has as points the
isotropic one-spaces, and as blocks their intersections with
secant projective lines. It is the classical
$S(2,q+1,q^3+1)$ design, with $q^2$ blocks through a point.
These design parameters and elementary finite-field Hermitian
linear algebra are the classical inputs used here.

Fix a block $L$. Its projective line corresponds to a nondegenerate
two-space $W$. A degenerate two-space would have a one-dimensional
radical and just one isotropic projective point, whereas a secant
has $q+1$. Write $V=W\perp Kz$ and use surjectivity of the field
norm to arrange $h(z,z)=1$. Every unital point outside $L$ has a
unique representative $v+z$, where
\[
                    v\in W,\qquad h(v,v)=-1.                \tag{H5}
\]
For distinct such vectors $v,w$, the projective line joining
$\langle v+z\rangle$ and $\langle w+z\rangle$ meets $W$ in
$\langle v-w\rangle$. Hence their joining block meets $L$ at
an isotropic point precisely when $v-w$ is isotropic. The point
of intersection specifies its direction, with no arbitrary
choice of a group isomorphism.

\paragraph{Corollary H.3 (full and deleted pencils).}
Let $\mathcal A$ be a set of at least two points on the unital
block $L$. Consider the bipartite geometry whose first part is
the unital points outside $L$, whose second part consists of
flags $(a,M)$ with $a\in\mathcal A$, $M\ne L$, and whose
incidence is $Q\in M$. This graph is connected for every prime
power $q\geq2$.

\emph{Proof.}
By (H5), its point-collinearity graph is exactly the shell graph
of Theorem H.2 for the directions represented by $\mathcal A$.
Every flag vertex has $q$ outside-point neighbors, so connectivity
of the point-collinearity graph connects the entire bipartite
graph. \hfill$\square$

Taking $\mathcal A=L$ gives the block residue of the Van Nypelseer
construction. Taking $\mathcal A=L\setminus\{p\}$ gives the
triangle residue in Part I, Theorem 4. Since $|L|=q+1$, both choices
retain two directions even at $q=2$. The proof thus justifies this
particular pencil deletion; it does not assume that arbitrary
deletion preserves connectivity.

\section{Formal scope and attribution}

The complete norm-shell proof was developed earlier in this project
for the conjectural statement in [VN, Section 4.6]. Its acceptance
and the question accounting are prior work; this edition does not
count that answer again. Formula (H2), the elementary unitary
calculation and the classical root generators are not claimed to
be historically new.

The retained module \texttt{HermitianPlaneTransport.lean} proves
thirteen statements: the Gram-equation equivalence, invariance of
the norm under the involution, norm preservation, existence and
uniqueness of (H2), and the root-matrix and nonzero-displacement
identities. They hold over an arbitrary field with an involution.
The retained \texttt{ResidueConnectivity.lean} supplies the abstract
generating-move connectivity implication. The finite-field
identification with $\operatorname{SL}(2,q)$, elementary generation,
projective map (H5) and the source geometry applications are
written proofs here. They are not formal consequences claimed to
have been instantiated in those modules.

\end{document}
